%
% Persönliche Macros
%
%

% Macros für Formeln
\newcommand{\jw}{j\omega}

% Begriffe

\newcommand{\OPV}{Operations\-ver\-stär\-ker}

% Unnummerierte Absatzüberschrift unterhalb der Section-Ebene.
% Der Titel erscheint ohne Gliederungsnummer auf der Subsection-Ebene
% des Inhaltsverzeichnisses. Die folgende Textzeile beginnt als eigener
% Absatz ohne Einzug.
\newcommand{\paragraphheading}[1]{%
    \par
    \addvspace{0.8\baselineskip}%
    \phantomsection
    \addcontentsline{toc}{subsection}{\protect\numberline{}#1}%
    \noindent\textbf{#1}\par\nobreak
    \addvspace{0.2\baselineskip}%
    \noindent\ignorespaces
}

% Gemeinsamer Stil für unveränderte Python-Auszüge im Codeanhang.
% Bei Bereichsauswahl Originalzeilen mit firstline/lastline/firstnumber setzen.
\lstdefinestyle{pythoncode}{
    language=Python,
    basicstyle=\ttfamily\small,
    keywordstyle=\bfseries,
    commentstyle=\color{black!60},
    stringstyle=\color{black!75},
    numbers=left,
    numberstyle=\ttfamily\scriptsize\color{black!60},
    numbersep=8pt,
    stepnumber=1,
    breaklines=true,
    breakatwhitespace=false,
    columns=fullflexible,
    keepspaces=true,
    showstringspaces=false,
    tabsize=4,
    captionpos=t,
    xleftmargin=2.5em,
    frame=single,
    rulecolor=\color{black!25},
    aboveskip=\baselineskip,
    belowskip=\baselineskip
}
\renewcommand{\lstlistingname}{Listing}
\captionsetup[lstlisting]{justification=raggedright,singlelinecheck=false}
